\chapter{Introducción}

\section{Contexto del proyecto}
\pending{Presentar la computación cuántica, el aprendizaje automático y el surgimiento del QML como contexto del trabajo.}

\section{Motivación}
\pending{Explicar la necesidad de una plataforma que simplifique la configuración, ejecución, repetición y comparación de experimentos QML.}

\section{Planteamiento del problema}
\pending{Definir el problema de construir un sistema de información que gestione el ciclo completo de los experimentos, desde los datos hasta la consulta de resultados.}

\section{Objetivos}
\subsection{Objetivo general}
\pending{Redactar el objetivo general de diseñar y desarrollar la plataforma modular de experimentación QML.}

\subsection{Objetivos específicos}
\pending{Enumerar objetivos relacionados con requisitos, arquitectura, motor QML, QCNN, interfaz, persistencia, reproducibilidad y extensibilidad.}

\section{Alcance}
\pending{Distinguir las funcionalidades implementadas, diseñadas, previstas y fuera del alcance de cada versión.}

\section{Metodología}
\pending{Describir el desarrollo incremental, la validación experimental, las pruebas automáticas y la gestión reproducible de configuraciones.}

\section{Estructura de la memoria}
\pending{Explicar brevemente el contenido y la relación entre los capítulos.}
